\section*{Appendix: Complete FSA-v5 Parameter Reference}
\addcontentsline{toc}{section}{Appendix: Complete FSA-v5 Parameter Reference}
\label{sec:appendix-v5-parameters}

This appendix is the canonical reference for every parameter in the
FSA-v5 model as implemented in the ``smc2fc``-port-ready codebase. For
each parameter we list its symbol, role in the equations of motion,
prior (if estimated) or frozen value, and where it first appears in the
main document.

The parameter set as implemented in
\path{models/fsa_high_res/estimation.py} consists of \textbf{37
estimated parameters} plus \textbf{14 frozen} (8 dynamics +
6 diffusion-related). The maximum-pinning rationale is the §11.6
recommendation derived from the Fisher-information analysis in §11.

\subsection*{Estimated parameters (37)}

These are the entries of \texttt{PARAM\_PRIOR\_CONFIG} in
\texttt{estimation.py}. Lognormal priors are parameterised as
$\log\theta \sim \mathcal{N}(\mu_{\log}, \sigma_{\log}^2)$ — the table
shows $(\exp(\mu_{\log}),\, \sigma_{\log})$. Normal priors are
parameterised directly as $(\mu, \sigma)$.

\begin{center}\small
\begin{tabular}{llll}
\toprule
Symbol & Role / location & Prior & Truth \\
\midrule
\multicolumn{4}{l}{\textit{Dynamics — aerobic capacity B}} \\
\midrule
$\tau_B$ & B time-scale / §11.1 & lognormal $(42.0,\,0.10)$ & $42.0$ d \\
$\kappa_B$ & B gain / §11.1 & lognormal $(0.01248,\,0.20)$ & $0.01248$ \\
$\epsilon_{AB}$ & autonomic-coupling on B / §11.1 & lognormal $(0.40,\,0.05)$ & $0.40$ \\
\midrule
\multicolumn{4}{l}{\textit{Dynamics — strength capacity S}} \\
\midrule
$\tau_S$ & S time-scale / §11.1 & lognormal $(60.0,\,0.10)$ & $60.0$ d \\
$\kappa_S$ & S gain / §11.1 & lognormal $(0.00816,\,0.20)$ & $0.00816$ \\
$\epsilon_{AS}$ & autonomic-coupling on S / §11.1 & lognormal $(0.20,\,0.05)$ & $0.20$ \\
\midrule
\multicolumn{4}{l}{\textit{Dynamics — fatigue F}} \\
\midrule
$\tau_F$ & F time-scale / §11.1 & lognormal $(6.36,\,0.15)$ & $6.36$ d \\
$\lambda_A$ & autonomic-coupling on F / §11.1 & lognormal $(1.00,\,0.05)$ & $1.00$ \\
\midrule
\multicolumn{4}{l}{\textit{Dynamics — variable-dose K (Busso 2003)}} \\
\midrule
$\mu_K$ & K damage rate / §11.1 & lognormal $(0.005,\,0.20)$ & $0.005$ \\
\midrule
\multicolumn{4}{l}{\textit{Dynamics — Stuart--Landau bifurcation parameter $\bar\mu$}} \\
\midrule
$\mu_0$ & baseline drive / §11.1 & lognormal $(0.036,\,0.20)$ & $0.036$ \\
$\mu_B$ & B$\to$A coupling / §11.1 & lognormal $(0.30,\,0.20)$ & $0.30$ \\
$\mu_S$ & S$\to$A coupling / §11.1 & lognormal $(0.15,\,0.20)$ & $0.15$ \\
$\mu_F$ & F$\to$A linear suppression / §11.1 & lognormal $(0.26,\,0.20)$ & $0.26$ \\
$\mu_{FF}$ & quadratic F penalty / §11.1 & lognormal $(0.40,\,0.05)$ & $0.40$ \\
$\eta$ & cubic damping in A / §11.1 & lognormal $(0.20,\,0.15)$ & $0.20$ \\
\midrule
\multicolumn{4}{l}{\textit{HR observation channel}} \\
\midrule
$\text{HR}_\text{base}$ & resting HR baseline (bpm) / §11.1 & normal $(62.0,\,2.0)$ & $62.0$ \\
$\kappa_B^{HR}$ & HR$\propto -B$ slope / §11.1 & lognormal $(12.0,\,0.15)$ & $12.0$ \\
$\alpha_A^{HR}$ & HR$\propto +A$ slope / §11.1 & lognormal $(3.0,\,0.20)$ & $3.0$ \\
$\beta_C^{HR}$ & HR circadian coefficient / §11.1 & normal $(-2.5,\,0.5)$ & $-2.5$ \\
$\sigma_{HR}$ & HR observation noise (bpm) / §11.1 & lognormal $(2.0,\,0.20)$ & $2.0$ \\
\midrule
\multicolumn{4}{l}{\textit{Sleep Bernoulli channel}} \\
\midrule
$k_C$ & sleep$\propto C$ slope / §11.1 & lognormal $(3.0,\,0.15)$ & $3.0$ \\
$k_A$ & sleep$\propto A$ slope / §11.1 & lognormal $(2.0,\,0.25)$ & $2.0$ \\
$\tilde c$ & sleep logit intercept / §11.1 & normal $(0.5,\,0.25)$ & $0.5$ \\
\midrule
\multicolumn{4}{l}{\textit{Stress observation channel}} \\
\midrule
$S_\text{base}$ & stress baseline / §11.1 & normal $(30.0,\,3.0)$ & $30.0$ \\
$k_F$ & stress$\propto +F$ slope / §11.1 & lognormal $(20.0,\,0.20)$ & $20.0$ \\
$k_{A,S}$ & stress$\propto -A$ slope / §11.1 & lognormal $(8.0,\,0.25)$ & $8.0$ \\
$\beta_C^{S}$ & stress circadian coefficient / §11.1 & normal $(-4.0,\,0.8)$ & $-4.0$ \\
$\sigma_S$ & stress observation noise / §11.1 & lognormal $(4.0,\,0.20)$ & $4.0$ \\
\midrule
\multicolumn{4}{l}{\textit{Steps observation channel (log-Gaussian)}} \\
\midrule
$\mu_\text{step,0}$ & log-steps intercept / §11.1 & normal $(5.5,\,0.3)$ & $5.5$ \\
$\beta_B^{st}$ & steps$\propto +B$ slope / §11.1 & lognormal $(0.8,\,0.20)$ & $0.8$ \\
$\beta_F^{st}$ & steps$\propto -F$ slope / §11.1 & lognormal $(0.5,\,0.25)$ & $0.5$ \\
$\beta_A^{st}$ & steps$\propto +A$ slope / §11.1 & lognormal $(0.3,\,0.25)$ & $0.3$ \\
$\beta_C^{st}$ & steps circadian coefficient / §11.1 & normal $(-0.8,\,0.2)$ & $-0.8$ \\
$\sigma_\text{st}$ & log-steps noise / §11.1 & lognormal $(0.5,\,0.15)$ & $0.5$ \\
\midrule
\multicolumn{4}{l}{\textit{Volume Load observation channel}} \\
\midrule
$\beta_S^{VL}$ & VL$\propto +S$ slope / §11.1 & lognormal $(100.0,\,0.15)$ & $100.0$ \\
$\beta_F^{VL}$ & VL$\propto -F$ slope / §11.1 & lognormal $(20.0,\,0.20)$ & $20.0$ \\
$\sigma_{VL}$ & VL observation noise / §11.1 & lognormal $(10.0,\,0.20)$ & $10.0$ \\
\bottomrule
\end{tabular}
\end{center}

\subsection*{Frozen parameters (8 dynamics + 6 diffusion-related)}

These do not appear in the estimated parameter vector. They are baked
into the ``EstimationModel.frozen\_params`` dict. Reasons for freezing
each one are tied directly to the §11 FIM analysis.

\begin{center}\small
\begin{tabular}{lllp{6cm}}
\toprule
Symbol & Frozen value & Role / location & Reason for freezing \\
\midrule
\multicolumn{4}{l}{\textit{Dynamics-side freezes}} \\
\midrule
$K_{FB}^0$  & $0.030$ & baseline aerobic K / §11.1   & Structurally non-identifiable. Slow-manifold relaxation collapses $K_{FB}^0$ and $\tau_K \mu_K$ into a single observable. §11.6 result (1). \\
$K_{FS}^0$  & $0.050$ & baseline strength K / §11.1  & Same as $K_{FB}^0$. \\
$\tau_K$    & $21.0$ d & K relaxation timescale / §11.1 & Busso-standard. Jointly weakly identified with $\mu_K$ (§11.6 result 5). \\
$n_\text{dec}$ & $4.0$ & Hill exponent / §10 eq.\ 25 & Structural shape parameter; not learned. \\
$B_\text{dec}$ & $0.07$ & B deconditioning threshold / §10 & Pinned to §10.4 Table 5 value. Inference requires $\ge 6$-week detraining episode (§11.6 result 3). \\
$S_\text{dec}$ & $0.07$ & S deconditioning threshold / §10 & Same. \\
$\mu_{B-}$ & $0.10$ & B deconditioning amplitude / §10 & Pinned to §10.4 Table 5 value. Inference requires same long detraining (§11.6 result 2). \\
$\mu_{S-}$ & $0.10$ & S deconditioning amplitude / §10 & Same. \\
\midrule
\multicolumn{4}{l}{\textit{Diffusion-side freezes (always frozen, both v4 and v5)}} \\
\midrule
$\sigma_B$ & $0.010$ & B Jacobi diffusion / §11.1 & Identifiability via diffusion is weak; standard practice. \\
$\sigma_S$ & $0.008$ & S Jacobi diffusion / §11.1 & Same. \\
$\sigma_F$ & $0.012$ & F CIR diffusion / §11.1 & Same. \\
$\sigma_A$ & $0.020$ & A CIR diffusion / §11.1 & Same. \\
$\sigma_K$ & $0.005$ & K CIR diffusion (shared) / §11.1 & Same. \\
$\phi$    & $0.0$    & circadian phase reference / §11.1 & Held at zero by convention. \\
\bottomrule
\end{tabular}
\end{center}

\subsection*{Initial state}

The 6D initial state is set from data or from \texttt{DEFAULT\_INIT}; it
has no explicit prior in the v5 \texttt{EstimationModel} (i.e.,
\texttt{INIT\_STATE\_PRIOR\_CONFIG = OrderedDict()}). Default values
correspond to a ``deconditioned but not unhealthy'' baseline:

\begin{center}\small
\begin{tabular}{lll}
\toprule
State & Default & Bounds \\
\midrule
$B$    & $0.05$  & $[0, 1]$ \\
$S$    & $0.10$  & $[0, 1]$ \\
$F$    & $0.30$  & $[0, \infty)$ \\
$A$    & $0.10$  & $[0, \infty)$ \\
$K_{FB}$ & $0.030$ & $[0, \infty)$ \\
$K_{FS}$ & $0.050$ & $[0, \infty)$ \\
\bottomrule
\end{tabular}
\end{center}

\subsection*{Reading guide for the smc2fc port}

After importing \texttt{models.fsa\_high\_res}, the port has access to:

\begin{itemize}
  \item \texttt{HIGH\_RES\_FSA\_V5\_ESTIMATION} — the ``EstimationModel`` to register with the SMC$^2$ engine. Built from the parameter prior config and frozen-params dict above.
  \item \texttt{HIGH\_RES\_FSA\_V5\_MODEL} — the ``SDEModel`` for forward simulation.
  \item \texttt{StepwisePlant} — the closed-loop plant for synthetic-feedback testing.
  \item \texttt{evaluate\_chance\_constrained\_cost} — the §9.6 v5 control cost (equations 23, 24).
\end{itemize}

The frozen values can be relaxed one at a time (move from
\texttt{frozen\_params} into \texttt{PARAM\_PRIOR\_CONFIG}) as the
calibration data set grows to support them.
